\documentclass[a4paper,10pt]{letter}
\usepackage[utf8]{inputenc}
\usepackage[T1]{fontenc}
\usepackage[italian]{babel}
\usepackage{lmodern}
\usepackage{geometry}
\geometry{
  top=1.8cm,
  bottom=2.0cm,
  left=2.5cm,
  right=2.5cm
  }
\usepackage{parskip}

\signature{Luca Conte} % Firma
\address{Treviso (TV)\\+39 338 6983681 \\luca.conte@gmail.com}
\date{\today}

\begin{document}

\begin{letter}{}

\opening{Gentili responsabili,}

mi chiamo Luca Conte, sono venuto a conoscenza della vostra realtà in occasione di un talk che avete tenuto presso Fabrica in provincia di Treviso (luogo in cui abito). Per me quella fu un'esperienza \textit{seminale}: costantemente pervaso da un effetto \textit{WoW-ma-allora-esistono-veramente-realtà-del-genere}, ho trovato molto interessanti i progetti presentati.  La qualità realizzativa percepita attraverso la presentazione e l'approccio tecnico/scientifico evinto dala visione dei vostri lavori, sono stati tutti fattori di sincera e profonda meraviglia. È stato piacevolissimo ritrovarmi nel workflow creativo che avete seguito nel realizzare \textit{Onirica}: pure io infatti nella realizzazione di \textit{Xynthesi}, mi ero ritrovato a far fronte agli stessi problemi generativi e, parlando con uno del vostro staff, ebbi modo di capire che gli step creativi da voi adottati erano stati sostanzialmente i medesimi adottati da me.

Ho un forte background tecnico scientifico che mi ha permesso di intraprendere una carriera soddisfacente nel mondo dell'IT, non ho mai comunque potuto fare a meno dell'arte in tutte le sue forme come complemento "spirituale" della mia persona sia in termini di attore passivo che attivo. Una decina di anni fa ho mosso i primi passi nel mondo della musica elettronica preferendo strumenti HW a quelli SW. Approdando nel corso della mia evoluzione artistica alla musica con synth modulari mi sono ritrovato a comporre suggestioni sonore che, per come la vedevo io (e per come la vedo tuttora), avrebbero potuto offrire esperienze più immersive  se accompagnati da dei visual audioreattivi. I primi passi li mossi con Processing, venendo da Java il passo era immediato, sapendo fin da subito che la mossa giusta sarebbe stata prima o poi passare a openFramewoks o in generale al C++. Un anno e mezzo fa con l'avvento degli LLM e di GPU con quantitativi di VRAM accettabili in ambito domestico mi sono avventurato nella sperimentazione di rendering di video musicali generati da IA (StableDiffusion su Deforum) pilotati tramite keyframing su feature sonore estratte (tramite Python) dai pezzi che componevo. Il risultato è stato il su citato \textit{Xynthesi} un viaggio audiovisivo di 16 minuti proiettato pubblicamente in apertura di una serata live di un camp modulare con un discreto successo. Terminato il progetto sono ritornato al creative coding implementando diversi algoritmi procedurali autoevolventi e sempre parametrizzati in tempo reale da input esterni. Questo ultimo step creativo è stata per me l'occasione di passare definitivamente a openFrameworks rafforzando e consolidando la mia conoscenza in C++.

  Faccio uso prevalentemente di strumenti Open Source e opero quasi esclusivamente su sistema operativo basato su Linux.  Ho avuto a che fare con varie architetture  oltre alla x86-64, tra cui ARM (SMT32/RaspberryPI) e AVR (Arduino). Amo il coding e l'eleganza della matematica che manifesta la sua estetica attraverso esso.

  Visitando il vostro sito, tornandoci dopo averlo visitato la prima volta a seguito della vostra uscita in Fabrica, per  vedere cosa stava bollendo in pentola, ho notato la posizione aperta e ho avvertito un forte impulso a candidarmi. Per pura trasparenza e in assoluta onestà penso che non avrei difficoltà a colmare eventuali hard skill necessarie per risultare produttivo, anche se certamente non potrei vendermi come senior nel contesto creativo che vi contraddistingue, avrei molto da imparare (e la cosa non mi spaventa), lato mio, tuttavia, il principale ostacolo che obbiettivamente vedo separarmi da quanto descritto nella candidatura sarebbe il full-time e il lavoro in presenza (non escludo che in un futuro la cosa possa cambiare), ma se siete aperti a una qualunque forma di collaborazione prevalentemente remota ma con possibili SAL \textit{de visu} da concordare ogni tot per me sarebbe un grosso privilegio, e costituirebbe una forte spinta motivazionale, entrare nella vostra realtà

Avrei rimpianto non averci provato, rimango a disposizione per un colloquio o una chiacchierata informale, qualora riteniate il mio profilo interessante.


\closing{Un cordiale saluto,}

\end{letter}

\end{document}
